\documentclass[conference]{IEEEtran}
\IEEEoverridecommandlockouts
\usepackage{cite}
\usepackage{amsmath,amssymb,amsfonts}
\usepackage{algorithmic}
\usepackage{graphicx}
\usepackage{textcomp}
\usepackage{xcolor}
\usepackage{booktabs}
\usepackage{microtype}
\usepackage{url}

\def\BibTeX{{\rm B\kern-.05em{\sc i\kern-.025em b}\kern-.08em
    T\kern-.1667em\lower.7ex\hbox{E}\kern-.125emX}}

\begin{document}

\title{PARISCV v2: Fast Microarchitectural Modeling and Custom Instruction Discovery for OpenHW RISC-V Cores}

\author{
    \IEEEauthorblockN{1\textsuperscript{st} Kushal R\IEEEauthorrefmark{1}, 2\textsuperscript{nd} Senior Author Name\IEEEauthorrefmark{2}, 3\textsuperscript{rd} Faculty Advisor\IEEEauthorrefmark{1}}
    \IEEEauthorblockA{\IEEEauthorrefmark{1}Department of Computer Science and Engineering, University Name, City, India}
    \IEEEauthorblockA{\IEEEauthorrefmark{2}Department of Electrical and Computer Engineering, Institution Name, City, India}
    \IEEEauthorblockA{Email: kushal@example.edu, senior@example.edu, advisor@example.edu}
}

\maketitle

\begin{abstract}
Evaluating Instruction Set Architecture (ISA) modifications and custom instruction offloading typically requires full Register-Transfer Level (RTL) simulation, which is too slow for early design-space exploration (DSE). Conversely, Instruction Set Simulators (ISS) like Spike run quickly but lack cycle-level timing fidelity. We present \textbf{PARISCV v2}, an automated pipeline for rapid custom-instruction discovery and ISA validation targeting OpenHW cores (CV32E40X with CV-X-IF and CV32E40P). By driving a calibrated 4-stage analytical microarchitectural cycle model directly from Spike execution commit traces, PARISCV v2 estimates real-core execution cycles in seconds without running RTL. Validated across 14 data points on the Embench-IoT suite against Verilator RTL ground truth, our cycle model achieves 0.00\%--3.45\% error. We demonstrate predicting the exact benefits of standard ISA extensions (Zbb, Xpulp) before synthesis and mining profitable custom instructions with tight prediction bounds. Furthermore, we reveal critical trade-offs between pipeline-native decoder edits and CV-X-IF coprocessor offload latencies, establishing the foundation for coarse-grained sub-word SIMD and TinyML acceleration.
\end{abstract}

\begin{IEEEkeywords}
RISC-V, OpenHW CV32E40X, CV-X-IF, Microarchitectural Modeling, Custom Instructions, TinyML.
\end{IEEEkeywords}

\section{Introduction}
Custom Instruction Set Extensions (ISE) are essential for meeting stringent energy, area, and throughput budgets in embedded edge systems. However, exploring custom instructions presents a major design dilemma:
\begin{itemize}
    \item \textbf{RTL Simulation (Verilator/FPGA)}: Provides exact ground-truth cycle counts but requires manual HDL modifications, testbench synthesis, and minutes to hours per run.
    \item \textbf{Instruction Set Simulators (ISS)}: Tools like Spike execute millions of instructions per second but omit microarchitectural hazards, branch penalties, and memory pipeline stalls.
\end{itemize}

Closest prior art, such as CIDRE~\cite{cidre2025}, analyzes application hotspots but lacks tightly coupled validation against standard open-source core interfaces like the CORE-V eXtension Interface (CV-X-IF). 

In this work, we present \textbf{PARISCV v2}, an end-to-end framework that bridges this gap. PARISCV v2 consumes standard Spike traces, reconstructs architectural register state, decomposes execution bottlenecks into microarchitectural stall categories, and scores custom candidate instructions in seconds. We validate predictions against real Verilator simulations of OpenHW CV32E40X and CV32E40P cores.

\section{System Architecture \& Methodology}

The PARISCV v2 architecture consists of three integrated components:

\subsection{Trace Ingestion \& State Reconstruction}
The target C program is compiled for the base ISA (\texttt{rv32imc}) and executed on Spike with commit logging enabled (\texttt{spike --log-commits}). The parser dynamically tracks architectural register writes, allowing exact divisor recovery for multi-cycle arithmetic (e.g., \texttt{div}/\texttt{rem}) to replace worst-case assumptions with exact formulaic costs ($3 + \text{CLZ}(\text{divisor})$).

\subsection{Microarchitectural Pipeline Model}
We model the documented 4-stage in-order pipeline of the CV32E40X/CV32E40P cores. Cycles are decomposed analytically:
\begin{equation}
\text{Cycles} = N_{\text{insn}} + \sum \text{Stalls}_{\text{hazard}} + \sum \text{Penalties}_{\text{control}}
\end{equation}
The model accounts for:
\begin{itemize}
    \item \textbf{Branch/Jump Penalties}: 3-cycle taken branch flush, 2-cycle jump flush, plus a $+1$ penalty for non-word-aligned 32-bit targets.
    \item \textbf{Data Hazards}: 1-cycle load-use penalties, \texttt{jalr} RAW register dependencies (1 cycle, or 2 cycles following a load).
    \item \textbf{Multi-cycle ALU}: Exact hardware loop execution, multiply-high latencies, and division leading-zero steps.
\end{itemize}

\subsection{Custom Instruction Candidate Mining}
The \texttt{find\_candidates} engine slides variable-width windows ($L \in [2, 4]$) across straight-line basic blocks in the trace. Candidates are normalized into canonical dataflow signatures. Constraints strictly reflect hardware feasibility:
\begin{enumerate}
    \item No control-flow divergence inside the fused sequence.
    \item Source register count $\le X\_NUM\_RS$ (typically 2 or 3).
    \item Exactly one live-out destination register ($X\_RFW\_WIDTH = 32$).
\end{enumerate}

\section{Experimental Validation \& Results}

\subsection{Cycle Model Accuracy vs. Verilator Ground Truth}
We evaluated PARISCV v2 across seven standard Embench-IoT kernels compiled for both CV32E40X and CV32E40P cores. Code layout was strictly pinned between Spike and RTL builds to eliminate base-address alignment variance.

\begin{table}[t]
\centering
\caption{PARISCV v2 Cycle Model Accuracy vs. Real Verilator RTL}
\label{tab:validation}
\resizebox{\columnwidth}{!}{%
\begin{tabular}{@{}lrrrrrr@{}}
\toprule
\textbf{Benchmark} & \textbf{Instrs} & \textbf{Model (40X)} & \textbf{RTL (40X)} & \textbf{40X Err} & \textbf{RTL (40P)} & \textbf{40P Err} \\ \midrule
\texttt{matmult-int} & 97,748    & 134,546    & 134,545    & \textbf{+0.00\%} & 134,547    & \textbf{-0.00\%} \\
\texttt{primecount}  & 2,273,871 & 3,927,267  & 3,927,224  & \textbf{+0.00\%} & 3,927,226  & \textbf{-0.00\%} \\
\texttt{edn}         & 49,365    & 65,190     & 65,179     & \textbf{+0.02\%} & 65,181     & \textbf{+0.01\%} \\
\texttt{tarfind}     & 53,846    & 118,587    & 117,149    & \textbf{+1.23\%} & 117,151    & \textbf{+1.23\%} \\
\texttt{md5sum}      & 52,607    & 72,497     & 71,400     & \textbf{+1.54\%} & 71,402     & \textbf{+1.53\%} \\
\texttt{statemate}   & 1,159     & 1,639      & 1,606      & \textbf{+2.05\%} & 1,608      & \textbf{+1.93\%} \\
\texttt{crc32}       & 24,608    & 30,764     & 29,737     & \textbf{+3.45\%} & 29,739     & \textbf{+3.45\%} \\ \bottomrule
\end{tabular}%
}
\end{table}

As shown in Table~\ref{tab:validation}, the model bounds execution error between \textbf{0.00\% and 3.45\%} across all 14 evaluation points. Stalls analysis indicates that control flow accounts for 8\%--25\% of all cycles, while target misalignment accounts for 3.6\%--9.1\% of cycles.

\subsection{Predicting Hardware Modifications Prior to RTL Build}
We tested predicting ISA modifications from a baseline trace before compilation or RTL synthesis:
\begin{itemize}
    \item \textbf{Zbb Extension (CV32E40X)}: Predicted 3,570 cycles saved on \texttt{edn}; RTL measured 3,669 cycles saved (\textbf{--2.7\% prediction error}).
    \item \textbf{Xpulp Hardware Loops (CV32E40P)}: Predicted a \textbf{3.243$\times$} speedup on single-loop MAC; RTL measured \textbf{3.241$\times$} (\textbf{--0.10\% error}).
\end{itemize}

\subsection{Fine-Grained vs. Coarse-Grained Coprocessor Insights}
To evaluate hardware offloading, we synthesized a CV-X-IF coprocessor implementing \texttt{pg.mac} ($r_d \leftarrow r_3 + r_1 \times r_2$) on CV32E40X. RTL simulation measured \textbf{zero speedup} (13,321 cycles baseline vs. 13,321 cycles with coprocessor). Because CV-X-IF requires a 2-cycle offload and writeback handshake, fusing 2-cycle scalar ALU operations is break-even. 

Conversely, native pipeline decoder edits (\texttt{pg.idx}, \texttt{pg.rol}) delivered up to \textbf{2,048 cycles saved} without offload overhead. This proves that \textbf{CV-X-IF is structurally unsuited for fine-grained scalar fusion and requires coarse-grained acceleration}.

\section{Ongoing \& Future Work: TinyML Acceleration}
Building on our finding that CV-X-IF requires multi-operation coarse-grained offloads, our current extension targets quantized edge AI. We are extending PARISCV v2 to detect packed sub-word SIMD opportunities (e.g., 4$\times$INT8 dot products $\sum_{i=0}^3 A_i B_i$). While standard RV32I requires 16+ instructions to unpack and compute a 4-byte dot product, our \texttt{pg\_xif\_tinyml} coprocessor executes the operation in a 2-cycle parallel hardware multiplier-adder tree, projecting a \textbf{3.5$\times$--5$\times$ throughput improvement} on INT8 Conv2D and GEMM inference kernels.

\section{Conclusion}
PARISCV v2 enables rapid, cycle-accurate design-space exploration for OpenHW RISC-V cores in seconds. Validated within 0.00\%--3.45\% of real RTL, it accurately predicts the benefits of ISA extensions and custom instructions. The framework provides a verified foundation for next-generation automated coprocessor synthesis and TinyML edge accelerators.

\begin{thebibliography}{00}
\bibitem{cidre2025} A. Author et al., ``CIDRE: Custom Instruction Discovery and Evaluation for RISC-V Processors,'' \emph{arXiv:2509.15782}, 2025.
\bibitem{openhw_cv32e40x} OpenHW Group, ``CV32E40X User Manual,'' \url{https://github.com/openhwgroup/cv32e40x}, 2024.
\bibitem{openhw_xif} OpenHW Group, ``CORE-V eXtension Interface (CV-X-IF) Specification,'' 2023.
\bibitem{embench} J. Bennett et al., ``Embench: An Open-Source Benchmark Suite for Embedded Systems,'' \emph{Free and Open Source Silicon Foundation}, 2019.
\end{thebibliography}

\end{document}
